% 图契约：同轴几何与特性阻抗；几何、方向和公式与正文同一约定。
\documentclass[tikz,border=5pt]{standalone}
\usepackage{fontspec,xeCJK,amsmath,amssymb}
\IfFontExistsTF{Times New Roman}{\setmainfont{Times New Roman}}{\setmainfont{TeX Gyre Termes}}
\IfFontExistsTF{Heiti SC}{\setCJKmainfont[AutoFakeBold=1.5]{Heiti SC}}{\setCJKmainfont[AutoFakeBold=1.5]{FandolHei-Regular}}
\usetikzlibrary{arrows.meta,calc,patterns}
\definecolor{ink}{HTML}{243745}
\definecolor{blue}{HTML}{275DAB}
\definecolor{green}{HTML}{177C63}
\definecolor{red}{HTML}{BD4935}
\definecolor{orange}{HTML}{AA6C16}
\definecolor{purple}{HTML}{715393}
\begin{document}
\begin{tikzpicture}[x=1cm,y=1cm,font=\fontsize{9.5}{12}\selectfont,
v/.style={-{Latex[length=2mm,width=1.3mm]},line width=.85pt},
axis/.style={-{Latex[length=1.7mm,width=1.1mm]},ink!65,line width=.55pt},
aux/.style={ink!45,dashed,line width=.55pt},
lab/.style={fill=white,inner sep=1.8pt},
card/.style={draw=ink!25,fill=blue!3,rounded corners=2pt,text width=6.1cm,minimum height=.85cm,inner sep=6pt,align=center}]
\path[use as bounding box](-8,-3.6)rectangle(8,3.65);
\node[font=\bfseries\fontsize{10.5}{13}\selectfont,text width=15.1cm,align=center]at(0.00000,3.22000){M7-049　同轴几何与特性阻抗};
\begin{scope}[shift={(-4,0)}]
\draw[ink,fill=blue!5](0,0)circle(2.05);
\draw[ink,fill=white](0,0)circle(0.7);
\fill[orange!15,even odd rule](0,0)circle(1.42)(0,0)circle(1.25);\draw[orange!70,dashed](0,0)circle(1.25);\draw[orange!70,dashed](0,0)circle(1.42);
\draw[aux](0.00000,0.00000)--(2.05000,0.00000);
\node[]at(0.70000,-0.30000){$a$};
\node[]at(2.15000,-0.30000){$b$};
\draw[v,blue](1.20000,0.60000)--(1.95000,0.97500);
\node[text=blue]at(2.20000,1.25000){$\mathbf E$};
\draw[v,green](1.20000,0.60000)--(0.75000,1.50000);
\node[text=green]at(-0.25000,1.55000){$\mathbf H$};
\draw[v,orange](0.00000,0.00000)--(0.88400,-0.88400);
\node[text=orange]at(0.35000,-0.75000){$r$};
\draw[orange,line width=2pt](0.88400,-0.88400)--(1.00400,-1.00400);
\draw[orange](1.00400,-1.00400)--(1.65000,-1.50000);
\node[text=orange]at(1.75000,-1.60000){$dr$};
\draw[red](-1.3,-0.55)circle(0.13);
\fill[red](-1.3,-0.55)circle(.035);
\node[text=red]at(-2.15000,-1.65000){$\mathbf S,+z$};
\node[]at(0.00000,-2.60000){环带面积 $2\pi r\,dr$};
\end{scope}
\node[card]at(4.00000,1.25000){$L'=\frac{\mu_0}{2\pi}\ln\frac ba,\quad C'=\frac{2\pi\varepsilon_0}{\ln(b/a)}$};
\node[card]at(4.00000,-1.05000){$Z_0=60\ln(2.30)\,\Omega\simeq50\,\Omega$};
\node[text width=15.2cm,align=center]at(0.00000,-3.20000){真空同轴线；忽略导体和介质损耗。};
\end{tikzpicture}\end{document}
